\capitulo{1}{Introducción}

Blender es un programa gratis que sirve para hacer objetos tridimensionales en 3D. Se usa mucho en diferentes campos, como para videojuegos, animaciones, para diseñar edificios o para la impresión 3D.

Cuando alguien está trabajando en Blender realiza muchas acciones, mueve objetos, cambia formas, usa distintas herramientas o ajusta materiales. Pero el programa no tiene ninguna función para ver automáticamente cómo trabaja la persona mientras modela.

Por eso, es difícil saber, por ejemplo:

\begin{itemize}
\item Cuáles son las herramientas que más se utilizan.
\item Cómo va cambiando la escena a medida que se modela.
\item Qué operaciones se realizan más a menudo.
\item Cómo se transforma la geometría del modelo con el tiempo.
\end{itemize}

Para solucionar este problema, en este trabajo final de carrera hemos creado un sistema que funciona con \textit{addons} de Blender. Este sistema graba automáticamente lo que se hace al modelar y luego nos deja analizar esos datos.

El sistema tiene dos partes clave:

\begin{itemize}
\item Un \textit{addon} que se encarga de guardar la información mientras se modela.
\item Otro \textit{addon} que sirve para analizar y mostrar esos datos de forma visual.
\end{itemize}

La idea principal de este proyecto es que sea más fácil entender cómo se realiza el proceso de modelado en Blender, porque el sistema recoge los datos automáticamente.

\section{Estructura de este trabajo}

Este trabajo está organizado en varios capítulos, donde explicamos todo el desarrollo del proyecto.

\begin{itemize}

\item \hyperref[cap:introduccion]{\textbf{Capítulo 1. Introducción}}\\
En el Capítulo 1, la Introducción, contamos de qué va el trabajo, qué problema queremos resolver y por qué nos interesó.

\item \hyperref[cap:objetivos]{\textbf{Capítulo 2. Objetivos del proyecto}}\\
En el Capítulo 2, los Objetivos del proyecto, hablamos de lo que queremos lograr con este sistema, tanto de forma general como más específica.

\item \hyperref[cap:conceptos]{\textbf{Capítulo 3. Conceptos teóricos}}\\
En el Capítulo 3, los Conceptos teóricos, explicamos algunas ideas importantes para entender bien el proyecto, como qué es Blender, los \textit{addons}, el análisis de datos o el formato CSV.

\item \hyperref[cap:tecnicas]{\textbf{Capítulo 4. Técnicas y herramientas}}\\
En el Capítulo 4, las Técnicas y herramientas, contamos qué tecnologías hemos usado, cómo está hecho el sistema por dentro y qué herramientas nos ayudaron a crearlo.

\item \hyperref[cap:desarrollo]{\textbf{Capítulo 5. Aspectos relevantes del desarrollo del proyecto}}\\
En el Capítulo 5, los Aspectos relevantes del desarrollo del proyecto, hablamos de las decisiones que tomamos al diseñar, los problemas que nos surgieron y cómo los resolvimos mientras hacíamos el sistema.

\item \hyperref[cap:relacionados]{\textbf{Capítulo 6. Trabajos relacionados}}\\
En el Capítulo 6, los Trabajos relacionados, mostramos otros trabajos e investigaciones que tienen que ver con estudiar cómo interactúa la gente en entornos 3D.

\item \hyperref[cap:conclusiones]{\textbf{Capítulo 7. Conclusiones y líneas de trabajo futuras}}\\
En el Capítulo 7, las Conclusiones y líneas de trabajo futuras, resumimos a qué conclusiones llegamos y proponemos qué se podría mejorar o hacer en el futuro.

\item \hyperref[cap:anexos]{\textbf{Anexos}}\\
En los Anexos hay más información que completa el proyecto.

\end{itemize}